\subsection{酉表示的希尔伯特直和与张量积}

设 G 是拓扑群。设 \((\varrho_{i})_{i\in\mathrm{I}}\) 是 G 在希尔伯特空间 \(E_{i}\) 中的一族酉表示。设 E 是各 \(E_{i}\) 的外部希尔伯特直和空间。对于 G 中每个 g 以及 E 中每个 \(x=(x_{i})_{i\in\mathrm{I}}\)，有

\[
\sum_ {i} \| \varrho_ {i} (g) x _ {i} \| ^ {2} = \sum_ {i} \| x _ {i} \| ^ {2} = \| x \| ^ {2},
\]

这证明了 \((\varrho_i(g)x_i)_{i\in \mathrm{I}}\) 属于 E 且与 \(x\) 有相同的范数。因此，元素 \(\varrho (g)\colon (x_{i})_{i\in \mathrm{I}}\mapsto (\varrho_{i}(g)x_{i})_{i\in \mathrm{I}}\) 是 \(\mathcal{L}(\mathrm{E})\) 中的酉元素。

引理 7. --- 映射 \(g \mapsto \varrho(g)\) 是 G 在 E 中的酉表示。

由 V, p. 380 的引理 4，只需证明对于 I 中每个 i 及 \(E_{i}\) 中每个 x，映射 \(g \mapsto \varrho_{i}(g)x\) 在 G 的单位元 e 处连续。该映射是连续映射 \(g \mapsto \varrho_{i}(g)x\) 与 \(E_{i}\) 到 E 的规范嵌入的复合，因而连续。

表示 \(\varrho\) 称为酉表示 \((\varrho_{i})_{i\in I}\) 的希尔伯特直和，记为 \(\varrho = \bigoplus \varrho_{i}\)。

设 G 和 H 是拓扑群。设 \(\varrho_{1}\)（分别为 \(\varrho_{2}\)）是 G（分别为 H）在希尔伯特空间 E \(_{1}\)（分别为 E \(_{2}\)）中的酉表示。设 E = E \(_{1}\) \(\widehat{\otimes}_{2}\) E \(_{2}\) 是 E \(_{1}\) 与 E \(_{2}\) 的希尔伯特张量积空间（EVT, V, p. 28, 定义 1）。对于 \((g,h)\in\mathrm{G}\times\mathrm{H}\)，令 \(\varrho(g,h)\) 为 E 的连续自同态 \(\varrho_{1}(g)\widehat{\otimes}_{2}\varrho_{2}(h)\)（TVS, V, p. 28）；无歧义时仍简记为 \(\varrho_{1}(g)\otimes\varrho_{2}(h)\)。

引理 8. --- a) 映射 \(\varrho\colon(g,h)\mapsto\varrho(g,h)\) 是 G × H 在 E 中的酉表示；

\begin{enumerate}
\def\labelenumi{\alph{enumi})}
\setcounter{enumi}{1}
\item
  对于任意标准正交基 \((e_i)_{i\in \mathrm{I}}\)，它是 \(\mathrm{E}_1\) 的基，从希尔伯特直和 \(\bigoplus_{i\in \mathrm{I}}\mathrm{E}_2\) 到 E 的映射 \((y_i)_{i\in \mathrm{I}}\mapsto \sum_{i\in \mathrm{I}}e_i\otimes y_i\) 是希尔伯特直和 \(\bigoplus_{i\in \mathrm{I}}\varrho_2\) 在 \(\operatorname {Res}_{\{e\} \times \mathrm{H}}^{\mathrm{G}\times \mathrm{H}}(\varrho)\) 中的等距 H-同构；
\item
  对于任意标准正交基 \((f_j)_{j\in \mathrm{J}}\)，它是 \(\mathrm{E}_2\) 的基，从希尔伯特直和 \(\bigoplus_{j\in \mathrm{J}}\mathrm{E}_1\) 到 E 的映射 \((x_{j})_{j\in \mathrm{J}}\mapsto \sum_{j\in \mathrm{J}}x_{j}\otimes f_{j}\) 是希尔伯特直和 \(\bigoplus_{j\in \mathrm{J}}\varrho_1\) 在 \(\operatorname {Res}_{\mathbf{G}\times \{e\}}^{\mathbf{G}\times \mathbf{H}}(\varrho)\) 中的等距 G-同构。
\end{enumerate}

映射 \((g,h)\mapsto\varrho(g,h)\) 是从 \(\mathrm{G}\times\mathrm{H}\) 到 \(\mathbf{GL}(\mathrm{E})\) 的同态（参见 EVT, V, p. 28, no. 2）。设 \((e_{i})_{i\in\mathrm{I}}\) 是 \(\mathrm{E}_{1}\) 的标准正交基。由 EVT, V, p. 29, prop. 3 和 cor. 2，映射

\[
u \colon (y _ {i}) \mapsto \sum_ {i \in I} e _ {i} \otimes y _ {i}
\]

是从希尔伯特直和 \(F_{2} = \bigoplus_{i \in I} E_{2}\) 到 E 上的等距映射。借助此等距映射，H 在 E 中的表示 \(\varrho_{\mathrm{H}} \colon h \mapsto \varrho(e, h)\) 被认同为表示族 \((\varrho_{2})_{i\in I}\) 在 \(F_{2}\) 中的直和表示。特别地，它是 H 在 E 中的酉表示（引理 7）。同样，同态 \(\varrho_{G}: g \mapsto \varrho(g,e)\) 是 G 在 E 中的酉表示。

取 \((g,h)\in\mathrm{G}\times\mathrm{H}\)。有 \(\varrho(g,h)=\varrho_{\mathrm{G}}(g)\circ\varrho_{\mathrm{H}}(h)\)，故 \(\varrho(g,h)\) 是酉的；此外，\(\varrho_{\mathrm{G}}(g)\) 与 \(\varrho_{\mathrm{H}}(h)\) 交换，所以由 V, p. 377 的引理 1，\(\varrho\) 是 \(G\times H\) 的酉表示。

最后，关于 \(\varrho\) 限制到子群 \(\{e\} \times \mathrm{H}\) 和 \(\mathrm{G} \times \{e\}\) 的断言已在上述论证过程中得到。

G×H 的表示 \((g,h)\mapsto\varrho_{1}(g)\otimes\varrho_{2}(h)\) 称为酉表示 \(\varrho_{1}\) 与 \(\varrho_{2}\) 的外张量积，记为 \(\varrho_{1}\boxtimes\varrho_{2}\)。

设 \(n \in \mathbf{N}\)，且 \((\mathrm{G}_i)_{1 \leqslant i \leqslant n}\) 是一族有限个拓扑群。设 \(\varrho_i\) 是 \(\mathrm{G}_i\) 在希尔伯特空间 \(\mathrm{E}_i\) 中的酉表示，对 \(1 \leqslant i \leqslant n\) 成立。类似地定义表示

\[
\varrho_ {1} \boxtimes \dots \boxtimes \varrho_ {n}
\]

\(G_{1} \times \cdots \times G_{n}\) 在希尔伯特空间 \(E = E_{1} \widehat{\otimes}_{2} \cdots \widehat{\otimes}_{2} E_{n}\) 中的表示（EVT, V, p. 27）。

若 \(G_{i} = G\)（\(1 \leqslant i \leqslant n\)），记 \(\Delta_{n}: G \to G^{n}\) 为由 \(g \mapsto (g, \ldots, g)\)（对所有 \(g \in G\)）定义的同态。我们用 \(\varrho_{1} \otimes \cdots \otimes \varrho_{n}\) 表示 G 的酉表示 \((\varrho_{1} \boxtimes \cdots \boxtimes \varrho_{n}) \circ \Delta_{n}\)。这称为酉表示 \(\varrho_{i}\) 的张量积。

对于每个置换 \(\sigma\)，即 \(\{1,\ldots,n\}\) 上的置换，规范同构

\[
\mathrm{E} _ {1} \widehat {\otimes} _ {2} \dots \widehat {\otimes} _ {2} \mathrm{E} _ {n} \rightarrow \mathrm{E} _ {\sigma (1)} \widehat {\otimes} _ {2} \dots \widehat {\otimes} _ {2} \mathrm{E} _ {\sigma (n)}
\]

（EVT, V, p. 28）是等距同构

\[
\varrho_ {1} \otimes \dots \otimes \varrho_ {n} \rightarrow \varrho_ {\sigma (1)} \otimes \dots \otimes \varrho_ {\sigma (n)}
\]

且是 G 的表示之间的同构。

若 \(\varrho_{i}=\varrho\) 对 \(1\leqslant i\leqslant n\) 成立，其中 \(\varrho\) 是 G 的酉表示，我们也用 \(\varrho^{\otimes n}\) 表示各 \(\varrho_{i}\) 的张量积表示，并称其为 \(n\) 次张量幂，即 \(\varrho\) 的张量幂。

